微网的日前购电计划必须在负荷和光伏尚未实现时确定，储能损耗、日内预报更新及紧急购电高价共同决定其运行费用。本文以十分钟为离散步长，建立统一的电力平衡、储能状态演化与购电结算模型，分别求解确定性日循环、历史预测下的固定计划、允许日内调整的计划以及波动电价下的调控问题。

针对问题一，将给定预测视为确定输入，在日初、日末储电量均为6000\,kWh的约束下建立线性规划。通过构造性论证说明，允许免费不利用余电时存在不同时充放电的同费用解，并以独立混合整数规划复核。最优日购电量为59482.70\,kWh，费用为35126.95元，两种求解所得费用一致。

针对问题二至四，使用严格按时间可见的历史曲线和预测残差构造风险分位净负荷，以线性规划确定购电量，再通过实际负荷驱动的储能反馈保证不缺供。问题三与问题4-3采用历史光伏和附件预报的凸组合；权重和分位只以1月15日至31日费用选择，2月至12月参数冻结。两种日内调整场景均选中光伏预报权重0.5、风险分位0.65；无调整场景采用纯历史光伏和0.8分位。正式评价期334天内，固定电价下无调整、有调整方案的费用分别为1390.82万元、1331.58万元；波动电价下分别为1467.58万元、1404.49万元。

为区分交易机会与预报信息的作用，固定重规划时刻、负荷修正及当前光伏观测，仅改变新小时预报的可用性。全部日内新预报在该对照下对应费用改善1.72万元，18时新预报未显示可辨认的经济优势。对于增购是否可撤回的合约歧义，分别进行冻结路径重计费与逐笔合约重新优化；后者费用为1346.05万元、1420.00万元。

进一步构造同等信息条件下的价格感知反馈基准和跨月滚动选参实验。二者在本数据上的累计费用均未优于原固定主方案，因此保留原主结果。本文仅对问题一给出确定性最优性结论；其余策略属于经过物理、费用与时序核验的可行启发式。由于模型族在分析同年数据后扩展，滚动检验仍为回顾式再分析，不构成新年度独立泛化证据。
